%% DSA 8301 — Statistical Inference for Big Data
%% Elsevier-style preprint layout | Title page + body with figures
\documentclass[11pt]{article}

%% ── Core packages ───────────────────────────────────────────
\usepackage[utf8]{inputenc}
\usepackage[T1]{fontenc}
\usepackage{amsmath,amssymb}
\usepackage{graphicx}
\usepackage{booktabs}
\usepackage{array}
\usepackage{caption}
\usepackage{subcaption}
\usepackage{float}
\usepackage{geometry}
\usepackage[hidelinks]{hyperref}
\usepackage{microtype}
\usepackage{enumitem}
\usepackage{parskip}
\usepackage{xcolor}
\usepackage{titlesec}
\usepackage{multirow}
\usepackage{setspace}
\usepackage{abstract}

%% ── Title-page fonts & decoration ───────────────────────────
\usepackage{mathpazo}      % Palatino-style serif (substitute for newpxtext)
\usepackage[scaled=0.92]{helvet}  % clean sans-serif for labels (substitute for cabin)
\renewcommand{\familydefault}{\rmdefault}
\usepackage{tikz}
\usepackage{tcolorbox}
\tcbuselibrary{skins,breakable}

%% ── Brand colours ───────────────────────────────────────────
\definecolor{titleblue}{RGB}{20, 52, 95}
\definecolor{accentgold}{RGB}{166, 124, 45}
\definecolor{rulelight}{RGB}{195, 195, 195}
\definecolor{boxbg}{RGB}{247, 249, 252}

%% ── Page geometry (body pages) ──────────────────────────────
\geometry{top=1.9cm, bottom=2.1cm, left=2.2cm, right=2.2cm}

%% ── Section heading sizes ───────────────────────────────────
\titleformat{\section}{\normalsize\bfseries}{\thesection.}{0.5em}{}
\titleformat{\subsection}{\small\bfseries}{\thesubsection}{0.5em}{}
\titleformat{\subsubsection}{\small\itshape}{\thesubsubsection}{0.4em}{}
\titlespacing{\section}{0pt}{6pt}{3pt}
\titlespacing{\subsection}{0pt}{4pt}{2pt}
\titlespacing{\subsubsection}{0pt}{3pt}{1pt}

%% ── Caption styling ─────────────────────────────────────────
\captionsetup{font=footnotesize, labelfont=bf, skip=3pt}
\setlength{\abovecaptionskip}{3pt}
\setlength{\belowcaptionskip}{1pt}
\setlength{\textfloatsep}{7pt plus 2pt minus 2pt}
\setlength{\intextsep}{5pt plus 1pt minus 1pt}
\setlength{\floatsep}{5pt}

%% ── Paragraph spacing ───────────────────────────────────────
\setlength{\parskip}{3.2pt}
\setlength{\parindent}{0pt}

%% ── Custom column types ─────────────────────────────────────
\newcolumntype{L}[1]{>{\raggedright\arraybackslash}p{#1}}
\newcolumntype{R}[1]{>{\raggedleft\arraybackslash}p{#1}}
\newcolumntype{C}[1]{>{\centering\arraybackslash}p{#1}}

%% ── Decorative ornament command ─────────────────────────────
\newcommand{\ornamentrule}{%
  \begin{center}
    \begin{tikzpicture}
      \draw[accentgold, line width=0.5pt] (0,0) -- (3.2,0);
      \fill[accentgold] (3.5,0) circle (1.4pt);
      \fill[accentgold] (3.8,0) circle (2.2pt);
      \fill[accentgold] (4.1,0) circle (1.4pt);
      \draw[accentgold, line width=0.5pt] (4.4,0) -- (7.6,0);
    \end{tikzpicture}
  \end{center}}

\graphicspath{{figures/}}

\renewcommand{\abstractname}{\large\bfseries Abstract}

%% ════════════════════════════════════════════════════════════
\begin{document}
%% ════════════════════════════════════════════════════════════

%% ── STANDALONE TITLE PAGE ───────────────────────────────────
\begin{titlepage}
\newgeometry{top=0cm, bottom=0cm, left=0cm, right=0cm}

\vspace*{2.0cm}

\begin{center}
{\fontsize{9.5}{12}\selectfont\sffamily\color{titleblue!70}\MakeUppercase{%
Strathmore Institute of Mathematical Sciences, iLabAfrica Research Centre}}\\[4pt]
\end{center}

\vspace{1.6cm}

\begin{center}
{\fontsize{25}{32}\selectfont\bfseries\color{titleblue}
Measuring Compounded\\[5pt]
Housing Vulnerability\\[5pt]
in Kenya}

\vspace{0.65cm}
\ornamentrule
\vspace{0.65cm}

{\fontsize{13}{18}\selectfont\itshape\color{black!70}
A Statistical Inference Framework Using\\[4pt]
the 2023/24 Kenya Housing Survey}
\end{center}

\vspace{1.6cm}

\begin{center}
\begin{tikzpicture}
  \draw[rulelight, line width=0.4pt] (0,0) -- (12cm,0);
\end{tikzpicture}
\end{center}

\vspace{1.5cm}

\begin{center}
{\fontsize{17}{22}\selectfont\bfseries\color{titleblue}Valerie Jerono}\\[5pt]
{\fontsize{10.5}{14}\selectfont\color{black!65}RegNo.\ \textbf{222331}}\\[3pt]
{\fontsize{9.5}{12}\selectfont\itshape\color{black!55}valerie.jerono@strathmore.edu}
\end{center}

\vspace{1.5cm}

\begin{center}
\begin{tcolorbox}[
  enhanced,
  width=13cm,
  colback=boxbg,
  colframe=titleblue!18,
  arc=3pt,
  boxrule=0.6pt,
  left=14pt, right=14pt, top=9pt, bottom=9pt
]
\centering
{\fontsize{10}{15}\selectfont
\renewcommand{\arraystretch}{1.4}
\begin{tabular}{r@{\hspace{10pt}}l}
{\sffamily\color{titleblue}\textbf{Course}}
  & DSA 8301: Statistical Inference for Big Data \\
{\sffamily\color{titleblue}\textbf{Lecturer}}
  & Prof.\ Jacob Ongalla \\
{\sffamily\color{titleblue}\textbf{Institution}}
  & Strathmore University, Nairobi, Kenya \\
{\sffamily\color{titleblue}\textbf{Dataset}}
  & Kenya Housing Survey 2023/24 (KNBS) \\
{\sffamily\color{titleblue}\textbf{Submission Date}}
  & 24 June 2026 \\
\end{tabular}}
\end{tcolorbox}
\end{center}

\vfill

\begin{center}
{\fontsize{8.2}{11}\selectfont\color{black!70}
\textsc{Keywords:}\enspace
Housing Vulnerability
$\cdot$ Composite Index
$\cdot$ Kenya Housing Survey
$\cdot$ Fractional Logit
$\cdot$ Nonparametric Inference\\[2pt]
Urban-Rural Disparity
$\cdot$ Insurance Pricing
$\cdot$ Sub-Saharan Africa}
\end{center}

\vspace{1.4cm}

\restoregeometry
\end{titlepage}

%% ── TITLE, AUTHOR, ABSTRACT (article-class equivalent of frontmatter) ──


\begin{abstract}
\small
Kenya's two-million-unit housing deficit and 2.3\% formal insurance penetration rate share a
root cause: the absence of a household-level instrument capable of quantifying compound
exposure across financial, structural, tenure, hazard, and utility dimensions simultaneously.
This study constructs and statistically validates the \textbf{Housing Financial Vulnerability
Score (HFVS)}, a five-pillar composite index derived from the 2023/24 Kenya Housing Survey
(KHS), covering 21,347 households across all 47 counties. A seven-stage cleaning pipeline
resolved five documented coding bugs and nullified 20,095 sentinel-contaminated cells,
yielding a 64-column model-ready frame. The national mean HFVS is 0.389 (SD\,$=0.057$),
statistically below the neutral midpoint of 0.5 (one-sample $t(21{,}346)=-285.5$,
$p<0.001$, $d=1.95$), with the 95\% bootstrap CI for the median at $(0.384,\,0.386)$.
Urban households score higher than rural counterparts (two-sample $t=18.22$, $p<0.001$,
$d=0.251$; Mann--Whitney $U=64{,}301{,}067$, $p<0.001$, $r=0.124$). Financial Stress (D1,
mean\,$=0.539$) significantly exceeds Utility Deprivation (D5, mean\,$=0.464$) at the
household level (Wilcoxon $W=71{,}276{,}896$, $p<0.001$, rank-biserial $r=0.324$),
confirming that vulnerability channels are unequal across dimensions. A monotone education
gradient is established by one-way ANOVA ($F(2,21344)=197.6$, $p<0.001$, $\eta^2=0.018$)
and confirmed by Kruskal--Wallis ($H=387.4$, $p<0.001$, $\varepsilon^2=0.018$). Multiple
OLS on six demographic proxies explains only 6.3\% of HFVS variance ($R^2=0.063$),
demonstrating that housing vulnerability cannot be inferred from income and education alone
and requires a structured housing intake instrument.

\vspace{4pt}
\noindent\textbf{Keywords:} Housing vulnerability; Composite index; Kenya Housing Survey;
Nonparametric inference; Urban-rural disparity; Education gradient; Gender;
Sub-Saharan Africa
\end{abstract}

\vspace{6pt}

%% ════════════════════════════════════════════════════════════
\section{Introduction}
\label{sec:intro}
%% ════════════════════════════════════════════════════════════

Kenya's housing crisis is both quantitative and qualitative. The national deficit stands at
two million units, growing by 200,000 annually, and concentrates in informal settlements
where land tenure is undocumented, building materials are non-permanent, and basic services
are simultaneously absent~\cite{knbs2024khs,unhabitat2022housing}. The April 2024 floods
crystallised the cost of this compound exposure: 270 deaths and 200,000 displacements, with a
casualty geography that tracked precisely to households that were simultaneously
flood-adjacent, tenure-insecure, structurally deficient, and financially stressed , not
any single condition, but all four at once.

The data needed to identify such households existed in the Kenya Housing Survey before the
rains arrived. What was missing was a validated instrument to translate those measurements
into an operational risk score. Three major policy mechanisms , the Affordable Housing
Programme (AHP), the Insurance Regulatory Authority (IRA) premium schedule, and the Kenya
Mortgage Refinance Company (KMRC) , all operate without a household-level vulnerability
indicator, helping explain why formal insurance penetration stands at only
2.3\%~\cite{ira2024annual} and formal mortgage take-up among KHS households at just
1.13\%~\cite{mwangi2024kenya}. This supply-side measurement failure has been systematically
misattributed to demand-side reluctance~\cite{owusu2024african}.

\subsection{Motivation for the Dataset}

The Kenya Housing Survey 2023/24 was selected because it is the first nationally
representative survey with sufficient breadth (443 raw variables across 14 linked files) and
depth (21,347 households, all 47 counties) to support a genuinely multidimensional
vulnerability instrument, rather than a single-proxy income or consumption index. No other
publicly available Kenyan dataset links financial, tenure, hazard, structural, and utility
indicators at the household level with national coverage.

\subsection{Objectives and Research Questions}

This study constructs and statistically validates the \textbf{Housing Financial Vulnerability
Score (HFVS)}, a five-pillar composite index, and addresses three research questions:

\begin{enumerate}[leftmargin=1.4em, topsep=2pt, itemsep=1pt, parsep=0pt]
\item \textbf{RQ1 (Distributional):} Is the national HFVS distribution significantly
  different from the neutral midpoint (0.5), and what is its shape?
\item \textbf{RQ2 (Group differences):} Do HFVS scores differ significantly across
  urban/rural strata, and how do the five pillars contribute to that difference?
\item \textbf{RQ3 (Predictive validity):} Can the composite be closely approximated from
  proxy variables alone, enabling a short intake instrument in place of a full housing
  survey?
\end{enumerate}
An affirmative answer to RQ3 means vulnerability scoring becomes deliverable via a
ten-minute intake form, with direct implications for actuarial
pricing~\cite{abunyewah2022multidimensional} and social programme
targeting~\cite{kimani2023tenure}.

%% ════════════════════════════════════════════════════════════
\section{Dataset Description}
\label{sec:data}
%% ════════════════════════════════════════════════════════════

\subsection{Source and Structure}

The KHS 2023/24 is produced by the Kenya National Bureau of Statistics (KNBS) under a
stratified multi-stage cluster design covering all 47 counties~\cite{knbs2024khs}. Fourteen
unit-level Stata files (household, individual, dwelling, land parcel, county, mortgage, loan,
NEMA, water-service, real estate, financiers, institutional, project information, and housing
types) were merged around the household spine (21,347 rows, 392 base columns) via the common
key \texttt{interview\_\_key}, augmented by county-level aggregates joined on county code.
The master frame reached 21,347 rows $\times$ 443 columns before cleaning. Survey weights
range from 24.9 to 8,162 (ratio: $327.6\times$); a 56.7\% divergence between weighted and
unweighted means for \texttt{wsvc\_sewer\_conns} confirms that all policy-relevant figures
must use weighted statistics throughout.

\subsection{Variable Definitions and Data Types}

The 447-column renamed master frame comprises 227 binary, 88 ordinal, 87 continuous, 25
free-text, 12 categorical, 4 identifier, and 4 zero-variance columns. The five HFVS pillars
are constructed from 53 engineered features:

\begin{itemize}[leftmargin=1.3em, topsep=2pt, itemsep=1pt, parsep=0pt]
\item \textbf{D1 -- Financial Stress} (12 features): log expenditure, housing-cost burden
  ratio, cost-burden flag, loan access, asset score, insurance coverage.
\item \textbf{D2 -- Tenure Security} (9 features): renter status, land ownership, title
  deed, lease documentation, eviction risk, tenure satisfaction.
\item \textbf{D3 -- Physical Hazard} (8 features): flood-zone and landslide-zone flags,
  terrain steepness, proximity hazard count.
\item \textbf{D4 -- Dwelling Quality} (10 features): wall/roof/floor durability,
  overcrowding, room count, floor area, dwelling age, structural-defect count.
\item \textbf{D5 -- Utility Deprivation} (9 features): electricity access, improved water
  and sanitation, cooking fuel, water-collection distance, sewer connectivity.
\end{itemize}

\subsection{Data Cleaning and Preprocessing}

Table~\ref{tab:pipeline} summarises the seven-stage pipeline from the 447-column renamed
master frame to the 64-column model-ready frame. Stage PL-01 classified all columns by type.
Stage PL-02 classified missingness mechanisms: structural zeros (\texttt{lp\_*} columns for
non-land-owning households), MAR skip-logic (rental vs.\ ownership modules, confirmed by
$r=-0.876$ anti-correlation between rent and dwelling-value missingness), and MCAR
(stratified median imputation by county $\times$ urban/rural). Stage PL-03 dropped 170
columns and nullified 20,095 sentinel-coded cells ($-1$, 98/99, 998/999) before capping
physical outliers at the 99th percentile. Stages PL-05--PL-07 applied near-zero-variance
pruning (prevalence $<3\%$ or $>97\%$), pairwise correlation pruning ($|r|>0.85$), and
variance-inflation-factor (VIF) screening (threshold $=10$; maximum VIF in the final set
$=5.33$). Five coding bugs were identified and corrected: an urban/rural code inversion, a
hazard-column/sanitation-column mismatch, an erroneously continuous \texttt{is\_slum}
indicator, a tenure-variable naming collision, and an incorrect dwelling-age source column.

\begin{table}[H]
\centering\small
\caption{Cleaning pipeline: column reduction from master frame to model-ready frame.}
\label{tab:pipeline}
\setlength{\tabcolsep}{4pt}
\begin{tabular}{@{}C{0.7cm}L{3.3cm}R{1.1cm}R{1.1cm}L{5.6cm}@{}}
\toprule
Stage & Action & Cols in & Cols out & Key decision \\
\midrule
PL01 & Column taxonomy        & 447 & 447 & 4 zero-variance flagged; 25 free-text identified \\
PL02 & Missingness audit      & 447 & 447 & MAR / MCAR / Structural classified per column \\
PL03 & Targeted drops + fixes & 447 & 277 & 170 cols removed; 20,095 sentinel cells nullified; 5 bugs fixed \\
PL04 & Feature engineering    & 277 & 330 & 53 HFVS features built across 5 pillars \\
PL05 & EDA \& imputation      & 330 & 330 & Stratified median fill; P99 outlier caps \\
PL06 & NZV + correlation      & 330 & 277 & 4 NZV + 6 correlated columns dropped \\
PL07 & VIF screening          & 277 & \textbf{64} & 0 dropped; max VIF $= 5.33$ \\
\bottomrule
\end{tabular}
\end{table}

Table~\ref{tab:descriptive} presents weighted descriptive statistics for the principal
continuous variables in the model-ready frame.

\begin{table}[H]
\centering\small
\caption{Weighted descriptive statistics for primary continuous variables ($n=21{,}347$).}
\label{tab:descriptive}
\setlength{\tabcolsep}{3.5pt}
\begin{tabular}{@{}L{3.6cm}R{1.1cm}R{1.2cm}R{1.1cm}R{1.1cm}R{1.0cm}R{1.0cm}@{}}
\toprule
Variable & Mean & Median & SD & IQR & Skew & Kurt \\
\midrule
HFVS composite         & 0.389 & 0.385 & 0.057 & 0.073 & $+0.44$ & $+0.30$ \\
D1 Financial Stress    & 0.539 & 0.534 & 0.078 & 0.110 & $+0.46$ & $-0.33$ \\
D2 Tenure Security     & 0.425 & 0.417 & 0.150 & 0.220 & $+0.15$ & $-0.82$ \\
D3 Hazard Exposure     & 0.128 & 0.000 & 0.218 & 0.330 & $+2.10$ & $+3.87$ \\
D4 Physical Quality    & 0.390 & 0.388 & 0.077 & 0.110 & $-0.08$ & $-0.72$ \\
D5 Utility Deprivation & 0.464 & 0.465 & 0.150 & 0.260 & $-0.12$ & $-1.20$ \\
log(total expenditure) & 10.39 & 10.44 & 0.672 & 0.899 & $-0.06$ & $+0.41$ \\
Housing burden ratio   & 0.092 & 0.000 & 0.119 & 0.127 & $+2.58$ & $+8.97$ \\
Persons per room       & 2.01  & 1.67  & 1.62  & 1.50  & $+2.14$ & $+8.83$ \\
Dependency ratio       & 0.706 & 0.500 & 0.704 & 1.000 & $+1.83$ & $+4.67$ \\
\bottomrule
\end{tabular}
\end{table}

%% ════════════════════════════════════════════════════════════
\section{Methodology}
\label{sec:method}
%% ════════════════════════════════════════════════════════════

\subsection{HFVS Construction}

The HFVS is an equal-weight arithmetic mean of five normalised pillar scores:
\begin{equation}
\text{HFVS} = \frac{D_1 + D_2 + D_3 + D_4 + D_5}{5}, \qquad D_i \in [0,1]
\label{eq:hfvs}
\end{equation}
where higher values indicate greater vulnerability. Each pillar $D_i$ is constructed by
z-score standardising each constituent indicator (with sign convention $s_j\in\{+1,-1\}$ so
that higher values indicate greater vulnerability), summing within-pillar, then min-max
rescaling to $[0,1]$~\cite{osei2021composite}. This z-score approach corrects for effective
weight suppression of D3 (Hazard Exposure), which has few active binary inputs and would
contribute $<5\%$ under a raw min-max scheme. The empirical 75th-percentile threshold
(HFVS\,$\geq 0.424$) identifies the high-vulnerability subpopulation ($\approx5{,}337$
households). An anti-leakage assertion tags all formula-ancestor variables and excludes them
from the proxy feature set used in the regression analysis by construction.

\subsection{Exploratory Data Analysis}

Exploratory analysis combined univariate distributional diagnostics (histograms, skewness,
excess kurtosis), Q-Q plots against the theoretical normal, and a Spearman correlation matrix
spanning the composite, the five pillars, and six core covariates. At $n=21{,}347$, the
Shapiro--Wilk (SW) test is trivially significant for any real-world variable, so normality
verdicts combine a SW test on a fixed $n=2{,}000$ subsample with effect-size criteria
($|\text{skew}|<0.5$ and $|\text{excess kurtosis}|<1.0$). Under this criterion the HFVS
composite is classified as approximately normal (skew $=-0.05$, kurtosis $=-0.41$); all five
individual pillars and the covariate set are classified as non-normal.

\subsection{Parametric Methods}

\textbf{(T1) One-sample $t$-test} ($H_0\colon \mu_\text{HFVS}=0.5$, two-tailed): tests
whether the national mean differs from the neutral midpoint, justified by the CLT at
$n=21{,}347$.

\textbf{(T2) Two-sample $t$-test, urban vs.\ rural} ($H_0\colon \mu_u=\mu_r$, Levene's
test for variance homogeneity applied first): tests the urban/rural HFVS gap with Cohen's
$d$ as the effect-size measure.

\textbf{(T3) Two-sample $t$-test, female- vs.\ male-headed households} (one-tailed,
$H_0\colon \mu_f\leq\mu_m$): tests whether female-headed households carry higher mean HFVS.

\textbf{(T4) One-way ANOVA} ($H_0\colon \mu_{T0}=\mu_{T1}=\mu_{T2}$, Tukey HSD
post-hoc): tests HFVS differences across three education tiers (none, primary/secondary,
post-secondary).

\textbf{(T5) Simple OLS}: regresses D1 Financial Stress on log(total expenditure) to isolate
the income component of financial vulnerability; heteroscedasticity assessed via
Breusch--Pagan test.

\textbf{(T6) Multiple OLS}: regresses the composite HFVS on six demographic proxies
(housing burden ratio, household size, dependency ratio, urban status, education tier,
log-expenditure) to benchmark explained variance from demographic predictors alone.

\textbf{(T7) Confidence intervals} (95\%, CLT-based): for the mean of each pillar and the
composite, confirming statistical distinguishability from the neutral midpoint.

\textbf{(M1) Maximum likelihood estimation}: fits a Beta$(\hat\alpha,\hat\beta)$
distribution to the composite via MLE with bootstrap standard errors ($B=500$), exploiting
the natural $[0,1]$ support of HFVS.

\subsection{Nonparametric Methods}

\textbf{(N1) Mann--Whitney $U$ test}: compares HFVS distributions between urban
($n=11{,}900$) and rural ($n=9{,}447$) strata; nonparametric complement to T2.

\textbf{(N2) Kruskal--Wallis test} with Bonferroni-corrected pairwise Mann--Whitney
post-hoc: nonparametric complement to the ANOVA (T4) for education-tier comparisons.

\textbf{(N3) Spearman rank correlation}: $\rho$ between D1 Financial Stress and the
composite, testing conceptual independence of pillars.

\textbf{(N4) Kolmogorov--Smirnov one-sample normality test}: applied to all five pillars
and the composite (Bonferroni-adjusted $\alpha=0.0083$), formally justifying the
parametric/nonparametric dual-method design.

\textbf{(N5) Wilcoxon signed-rank test}: tests whether D1 and D5 share equal
household-level medians ($H_0\colon \theta_{D1-D5}=0$), validating dimension heterogeneity
without assuming normality of the paired differences~\cite{torabi2022utility}. Asymptotic
Relative Efficiency (ARE) of the Wilcoxon vs.\ the paired $t$-test is reported.

\textbf{(N6) Bootstrap confidence intervals} ($B=5{,}000$ resamples): for the national
mean, median, and 75th-percentile HFVS, providing non-parametric intervals where no
closed-form alternative exists~\cite{li2023bootstrap}.

\textit{Assumptions.} Parametric tests rely on the CLT at $n=21{,}347$. OLS additionally
assumes homoscedastic, uncorrelated errors; Breusch--Pagan tests are reported. Nonparametric
tests assume only independent observations and (for Wilcoxon) symmetric paired differences.

\textit{Python packages:} \texttt{pandas}, \texttt{numpy}, \texttt{scipy},
\texttt{statsmodels}, \texttt{scikit-learn}, \texttt{matplotlib}, \texttt{seaborn}.

%% ════════════════════════════════════════════════════════════
\section{Results}
\label{sec:results}
%% ════════════════════════════════════════════════════════════

\subsection{Distributional Profile and Normality}

Figure~\ref{fig:pillar_dist} shows that individual pillars are clearly non-normal: D3
(Hazard Exposure, mean\,$=0.128$) is severely right-skewed, with the majority of households
reporting zero flood-zone classification; D5 (Utility Deprivation, mean\,$=0.464$) is
bimodal. A Kolmogorov--Smirnov normality test (Bonferroni-adjusted $\alpha=0.0083$ for six
simultaneous tests) rejects normality for all five pillars and the composite
($p<0.001$ in every case), formally justifying the dual parametric/nonparametric design.
All parametric tests rely on the CLT, not on individual-observation normality.

\begin{figure}[h!]
  \centering
  \includegraphics[width=4\linewidth, height=6.4cm, keepaspectratio]{fig1_pillar_distributions.png}
  \caption{HFVS composite and five-pillar score distributions ($n=21{,}347$, weighted). The
  composite is near-symmetric; individual pillars are non-normal. D3 (Hazard) is severely
  right-skewed; D5 (Utility Deprivation) is bimodal.}
  \label{fig:pillar_dist}
\end{figure}

Figure~\ref{fig:qq} confirms this pattern: the HFVS Q-Q plot hugs the theoretical diagonal
through the central 95\% of the distribution, while expenditure, persons-per-room,
housing-burden ratio, and dependency ratio show pronounced departures from normality ,
justifying the mixed parametric/nonparametric inferential strategy adopted throughout.

\begin{figure}[h!]
  \centering
  \includegraphics[width=1.95\linewidth, height=5.4cm, keepaspectratio]{fig2_qq_plots.png}
  \caption{Q-Q plots for the HFVS composite and five representative covariates. The composite
  (top-left) closely follows the theoretical normal line; the remaining variables show clear
  departures, validating the use of nonparametric methods for individual-variable inference.}
  \label{fig:qq}
\end{figure}

\subsection{Parametric Results}

\textbf{T1 — One-sample $t$-test.} The national mean HFVS is 0.389 (SD\,$=0.057$, $n=21{,}347$).
Testing $H_0\colon\mu=0.5$ yields $t(21{,}346)=-285.5$, $p<0.001$, 95\% CI for the mean
$(0.388,\,0.390)$. Cohen's $d=1.95$ reflects the large distance from the midpoint in
standard-deviation units. The national mean sits well below 0.5, confirming that most
households occupy the lower-vulnerability half of the scale.

\textbf{T2 — Two-sample $t$-test, urban vs.\ rural.} Levene's test ($F=1.82$, $p=0.18$)
confirms equal variances. Urban households (mean\,$=0.395$, $n=11{,}900$) score
significantly higher than rural households (mean\,$=0.381$, $n=9{,}447$):
$t(21{,}345)=18.22$, $p<0.001$, Cohen's $d=0.251$ (small). Urban households face greater
financial stress (D1) and tenure insecurity (D2) from informal-settlement rent burdens;
rural households face greater utility deprivation (D5)~\cite{kimani2023tenure}.

\textbf{T3 — Two-sample $t$-test, female- vs.\ male-headed households.} Female-headed
households (mean\,$=0.392$, $n=6{,}817$) score higher than male-headed (mean\,$=0.388$,
$n=14{,}530$): $t=5.53$, $p<0.001$ (one-tailed), Cohen's $d=0.081$ (small). The gap is
statistically robust but modest in magnitude.

\textbf{T4 — One-way ANOVA, education tier.} Levene's test ($p=0.080$) confirms variance
homogeneity. Means: no education\,$=0.395$; primary/secondary\,$=0.383$;
post-secondary\,$=0.372$. $F(2,21344)=197.6$, $p<0.001$, $\eta^2=0.018$ (small). Tukey HSD
post-hoc confirms all pairwise differences ($p<0.001$): a monotone protective
gradient — more education, lower vulnerability — but education explains only 1.8\% of
HFVS variance, confirming that it does not insulate households from hazard or dwelling
quality deficits.

\textbf{T5 \& T6 — OLS regressions.} Simple OLS of D1 on log(expenditure) yields
$R^2=0.072$ ($\beta_1=-0.032$, $t=-40.7$, $p<0.001$; Breusch--Pagan $p<0.001$, indicating
heteroscedasticity). Multiple OLS regressing the composite on six demographic proxies
(housing burden ratio, household size, dependency ratio, urban status, education tier,
log-expenditure) yields $R^2=0.063$ ($F(6,21340)=238.5$, $p<0.001$; all VIF\,$<5$).
Demographic predictors jointly explain only 6.3\% of composite variance, with 93.7\%
captured by the engineered housing conditions measured in the full survey instrument.

\textbf{T7 — Confidence intervals.} The 95\% CLT-based CI for the composite mean is
$(0.388,\,0.390)$ (width\,$=0.0015$); no dimension CI includes 0.5, confirming all
pillar means are statistically distinguishable from the neutral midpoint.

\textbf{M1 — MLE Beta fit.} The composite HFVS\,$\in[0,1]$ is modelled as
Beta$(\hat\alpha,\hat\beta)$ via MLE (bootstrap SE, $B=500$). Estimates:
$\hat\alpha=28.67$ (SE\,$=0.31$), $\hat\beta=44.97$ (SE\,$=0.50$);
log-likelihood\,$=31{,}098$; AIC\,$=-62{,}192$; BIC\,$=-62{,}176$. The implied Beta mean
(0.3893) matches the empirical mean (0.3892) to four decimal places, confirming the Beta
family as a suitable parametric approximation for this bounded composite.

Table~\ref{tab:allresults} summarises all parametric and nonparametric test statistics.

\begin{table}[H]
\centering\small
\caption{Summary of all parametric and nonparametric statistical tests.}
\label{tab:allresults}
\setlength{\tabcolsep}{3pt}
\begin{tabular}{@{}L{2.4cm}L{3.4cm}R{2.3cm}R{0.9cm}L{3.1cm}@{}}
\toprule
Test & Hypothesis & Statistic & $p$ & Conclusion \\
\midrule
\multicolumn{5}{@{}l}{\textit{Parametric}} \\[1pt]
T1 One-sample $t$ & $H_0$: $\mu=0.5$ & $t=-285.5$ & $<.001$ & Mean 0.389; below midpoint \\
T2 Two-sample $t$ & $H_0$: Urban$=$Rural & $t=18.22$, $d=0.251$ & $<.001$ & Urban $>$ Rural \\
T3 Two-sample $t$ & $H_0$: Female$\leq$Male & $t=5.53$, $d=0.081$ & $<.001$ & Female $>$ Male \\
T4 One-way ANOVA & $H_0$: $\mu_{T0}=\mu_{T1}=\mu_{T2}$ & $F=197.6$, $\eta^2=0.018$ & $<.001$ & Monotone edu gradient \\[3pt]
T5 Simple OLS & $H_0$: $\beta_1=0$ & $R^2=0.072$ & $<.001$ & 7.2\% of D1 variance \\
T6 Multiple OLS & $H_0$: all $\beta_j=0$ & $R^2=0.063$ & $<.001$ & 6.3\% of HFVS variance \\
M1 Beta MLE & HFVS$\sim$Beta$(\alpha,\beta)$ & AIC$=-62{,}192$ & -- & $\hat\alpha=28.67$, $\hat\beta=44.97$ \\[3pt]
\multicolumn{5}{@{}l}{\textit{Nonparametric}} \\[1pt]
N1 Mann-Whitney & $H_0$: Urban$=$Rural & $U=64{,}301{,}067$, $r=0.124$ & $<.001$ & Urban $>$ Rural \\
N2 Kruskal-Wallis & $H_0$: edu tiers equal & $H=387.4$, $\varepsilon^2=0.018$ & $<.001$ & Monotone gradient \\
N3 Spearman & HFVS vs.\ D1 & $\rho=0.200$ & $<.001$ & Moderate positive \\
N4 KS normality & $H_0$: all dims Normal & KS D1$=0.042$, D3$=0.266$ & $<.001$ & All non-normal \\
N5 Wilcoxon & $H_0$: D1$-$D5 diff$=0$ & $W=71{,}276{,}896$, $r=0.324$ & $<.001$ & D1$>$D5; pillars differ \\
N6 Bootstrap & 95\% CI for mean & $[0.388,\,0.390]$ & -- & Confirms T1 \\
\bottomrule
\end{tabular}
\end{table}

\subsection{Nonparametric Results}

\textbf{N1 — Mann--Whitney $U$ test.} Urban households (median\,$=0.390$) score higher
than rural households (median\,$=0.377$): $U=64{,}301{,}067$, $Z=18.09$, $p<0.001$,
rank-biserial $r=0.124$ (small). The Hodges--Lehmann shift estimate is 0.013 on the
$[0,1]$ scale — statistically detectable at this sample size but practically modest.
This nonparametric result confirms T2~\cite{kimani2023tenure}.

\textbf{N2 — Kruskal--Wallis, education tier.} $H(2)=387.4$, $p<0.001$,
$\varepsilon^2=0.018$ (small). Bonferroni-corrected pairwise Mann--Whitney tests confirm
all three education-tier pairs differ ($p<0.001$), replicating the monotone gradient
found under ANOVA (T4).

\textbf{N3 — Spearman correlation.} D1 Financial Stress vs.\ the composite:
$\rho=0.200$ ($p<0.001$, 95\% CI\,$[0.188,\,0.213]$). The value $\rho<1$ confirms that
the remaining four dimensions contribute independent variance, validating the
multi-pillar design~\cite{osei2021composite}.

\textbf{N5 — Wilcoxon signed-rank test.} The test on paired household-level differences
$(D1-D5)$ yields $W=71{,}276{,}896$, $p<0.001$, rank-biserial $r=0.324$. The mean D1
(Financial Stress\,$=0.539$) exceeds the mean D5 (Utility Deprivation\,$=0.464$) by 0.076
units, and the median difference is 0.075. \emph{Financial Stress is systematically the
higher-scoring dimension at the household level}, not Utility Deprivation. The Asymptotic
Relative Efficiency of Wilcoxon vs.\ the paired $t$-test is $3/\pi\approx0.955$ under
normality; under the observed non-normal differences the Wilcoxon is at least as
efficient~\cite{torabi2022utility,osei2021composite}.

\textbf{N6 — Bootstrap confidence intervals.} Using $B=5{,}000$ resamples
(Figure~\ref{fig:spearman}), the 95\% bootstrap CI for the national mean is
$(0.388,\,0.390)$ and for the median is $(0.384,\,0.386)$. The 75th-percentile
(high-vulnerability threshold) CI is $(0.423,\,0.425)$, identifying the $\approx5{,}337$
households in the top quartile~\cite{li2023bootstrap}.

\subsection{County-Level Vulnerability Profile}

Figure~\ref{fig:counties} (left panel) ranks the 15 highest-vulnerability counties by
weighted mean HFVS. Mandera leads with a weighted mean of 0.556, followed by Trans Nzoia
(0.530) and Tana River (0.502), all materially above the national mean. The right panel of
Figure~\ref{fig:counties} relates county-level HFVS to the triple-exposure rate , the share
of households that are simultaneously flood-adjacent, tenure-insecure, and rent-stressed,
which is the exact profile implicated in the April 2024 flood casualties. Counties without a
formal mortgage market (shown in red) cluster systematically above HFVS $\approx0.42$ and
exhibit higher triple-exposure rates, evidence that financial exclusion and housing
vulnerability mutually reinforce each other~\cite{mwangi2024kenya,owusu2024african}.

\begin{figure}[h!]
  \centering
  \begin{subfigure}[t]{0.47\linewidth}
    \centering
    \includegraphics[width=\linewidth, height=5.0cm, keepaspectratio]{fig7_top_counties.png}
    \caption{Top 15 counties by weighted mean HFVS.}
    \label{fig:top_counties}
  \end{subfigure}
  \hfill
  \begin{subfigure}[t]{0.49\linewidth}
    \centering
    \includegraphics[width=\linewidth, height=5.0cm, keepaspectratio]{fig8_county_triple_vs_hfvs.png}
    \caption{County HFVS vs.\ triple-exposure rate.}
    \label{fig:county_triple}
  \end{subfigure}
  \caption{County-level vulnerability profile across all 47 counties. Mandera, Trans Nzoia,
  and Tana River lead by weighted mean HFVS (left). Mortgage-absent counties (red, right)
  cluster in the high-HFVS, high-triple-exposure quadrant, linking financial exclusion to
  compound exposure.}
  \label{fig:counties}
\end{figure}

%% ════════════════════════════════════════════════════════════
\section{Discussion}
\label{sec:discussion}
%% ════════════════════════════════════════════════════════════

\subsection{Interpretation of Key Findings}

The HFVS composite meets all design criteria for a usable policy instrument: it is stable
(bootstrap CI width $<0.002$), nationally representative, non-degenerate, and statistically
distinguishable from a neutral midpoint. The low inter-pillar Spearman correlations
confirm that each dimension captures a genuinely distinct vulnerability
channel~\cite{osei2021composite} — the core empirical justification for a five-pillar
rather than a single-proxy design.

The Wilcoxon result carries a direct policy implication: D1 Financial Stress (mean\,$=0.539$)
exceeds D5 Utility Deprivation (mean\,$=0.464$), meaning financial pressure is the dominant
unpriced exposure in Kenya's residential portfolio at the household level. Flat national
premium schedules that proxy vulnerability with utility access alone will systematically
misclassify the financially stressed subpopulation, a pattern documented across West African
cities~\cite{torabi2022utility,nkeki2023flood}. The urban--rural gap, while statistically
robust ($d=0.251$), is driven by distinct channel composition: urban households carry higher
D1 and D2 (tenure insecurity and rent burden); rural households carry higher D5 (utility
deprivation).

\subsection{Comparison of Parametric and Nonparametric Methods}

For the urban--rural comparison, T2 ($t=18.22$) and N1 (Mann--Whitney $U$, $r=0.124$) agree
exactly: urban households score higher, effect size is small. For education tiers, T4 (ANOVA
$F=197.6$, $\eta^2=0.018$) and N2 (Kruskal--Wallis $H=387.4$, $\varepsilon^2=0.018$) agree
identically. This convergence across all six substantive comparisons confirms findings are
robust and not artefacts of distributional assumptions. For non-normal variables
— individual pillars, expenditure, and burden ratios — nonparametric methods are the
only valid primary tools, since parametric assumptions are violated at the variable level.
The contrast between simple OLS ($R^2=0.072$ for D1) and multiple OLS ($R^2=0.063$ for
the composite) shows that demographic proxies explain very little of either the pillar or
the composite variance.

\subsection{Unexpected Results}

One result was not anticipated at the outset: urban households score higher than rural
households on the composite despite generally higher incomes. Dimension-level decomposition
shows this is driven by tenure insecurity (D2) and rent burden (D1) in informal settlements
outweighing rural disadvantages in utility access (D5). The education gradient
($\eta^2=0.018$) is also notable: highly significant but explaining only 1.8\% of variance,
confirming that education is not a sufficient proxy for housing vulnerability.

\subsection{Practical Implications}

For the \textbf{IRA}, county HFVS rankings provide a validated, evidence-based basis for
county-differentiated premium loading. For \textbf{AHP} site selection, the pillar
decomposition identifies counties where dwelling quality is the binding constraint ,
directing interventions toward structural upgrading , versus counties where hazard
exposure dominates, directing interventions toward relocation and flood
mitigation~\cite{abunyewah2022multidimensional,nkeki2023flood}. For \textbf{KMRC},
mortgage-absent counties in the upper HFVS quartile represent the segment where the social
return to housing-finance extension is highest.

%% ════════════════════════════════════════════════════════════
\section{Conclusions}
\label{sec:conclusion}
%% ════════════════════════════════════════════════════════════

This study constructed and statistically validated the Housing Financial Vulnerability
Score, a five-pillar composite index built from the 2023/24 Kenya Housing Survey. The
principal findings are:

\begin{enumerate}[leftmargin=1.5em, topsep=2pt, itemsep=2pt, parsep=0pt]
\item The national mean HFVS is 0.389 ($t=-285.5$, $p<0.001$, $d=1.95$); the bootstrap
  95\% CI for the mean is $(0.388,\,0.390)$, confirming the national average is
  significantly below the neutral midpoint of 0.5.
\item The five HFVS pillars are mutually near-orthogonal, validating the multi-dimensional
  construction. Spearman $\rho$ between D1 and the composite is only 0.200, ruling out
  dimension redundancy.
\item Financial Stress (D1, mean\,$=0.539$) is systematically \emph{higher} than Utility
  Deprivation (D5, mean\,$=0.464$) at the household level (Wilcoxon $W=71{,}276{,}896$,
  $r=0.324$, $p<0.001$); single-proxy utility indices would misclassify a substantial share
  of financially stressed households.
\item Urban households carry higher HFVS scores than rural households ($t=18.22$,
  $d=0.251$; Mann--Whitney $r=0.124$; both $p<0.001$), driven by distinct vulnerability
  channels in each stratum.
\item A monotone education gradient exists ($F=197.6$, $\eta^2=0.018$, $p<0.001$), but
  education explains only 1.8\% of HFVS variance, confirming vulnerability cannot be
  proxied from socioeconomic status alone.
\item Female-headed households score significantly higher than male-headed households
  ($t=5.53$, $d=0.081$, $p<0.001$).
\item Demographic proxy variables explain only 6.3\% of HFVS variance (multiple OLS,
  $R^2=0.063$), underscoring that housing vulnerability requires a structured housing
  intake instrument rather than demographic screening alone.
\end{enumerate}

\subsection{Limitations}

D3 (Hazard Exposure) is structurally underweighted because the full proximity sub-column
battery was absent from this data extract, limiting the pillar to flood- and landslide-zone
flags alone. The county-level supply-side features (mortgage presence, sewer connectivity)
are constants within each county, so they capture only between-county rather than
within-county variance. This is a single cross-sectional survey wave, so all relationships
reported here are correlational; no causal claim is made about, for example, whether urban
residence \emph{causes} higher HFVS. Finally, the demographic-only classification AUC of
approximately 0.55 confirms that HFVS cannot be replicated from administrative data alone,
which bounds how far the proxy-instrument approach in RQ3 can be pushed without at least some
structural housing variables.

\subsection{Recommendations for Future Work}

Future iterations should incorporate satellite-derived hazard layers to correct the D3
under-weighting, following recent East African remote-sensing
approaches~\cite{mugiraneza2022rwanda}. Longitudinal validation , linking HFVS scores to
actual insurance claims or housing-finance default records , would provide the definitive
test of predictive validity that a single cross-section cannot offer. Sub-county
(ward-level) spatial grouping would also strengthen the external validity of any
machine-learning proxy model built on top of the HFVS framework.

%% ════════════════════════════════════════════════════════════
\section*{Acknowledgements}
%% ════════════════════════════════════════════════════════════
\small
Microdata accessed via the KNBS Microdata Portal
(\url{https://statistics.knbs.or.ke/nada/index.php/catalog/184}).

%% ════════════════════════════════════════════════════════════
\bibliographystyle{unsrt}
\bibliography{references}
%% ════════════════════════════════════════════════════════════

\end{document}